\section{凸集的典型例子}

\subsection{超平面与半空间}

\begin{BoxDefinition}[超平面]
    超平面（Hyperplane）定义为
    \begin{Equation}[超平面]
        \qty{\vb{x}\mid\vb{a}^T\vb{x}=b}\qquad \vb{a}\neq \vb{0}
    \end{Equation}
\end{BoxDefinition}

\begin{BoxDefinition}[半空间]
    半空间（Hyperplane）定义为
    \begin{Equation}[半空间]
        \qty{\vb{x}\mid\vb{a}^T\vb{x}\leq b}\qquad \vb{a}\neq \vb{0}
    \end{Equation}
\end{BoxDefinition}

如\cref{fig:超平面,fig:半空间}所示，超平面会将空间划分为两个部分，半空间就是其中一侧的空间。

\begin{Figure}[超平面与半空间]
    \figuresub[超平面]{\includegraphics{Hyperplane.fig.pdf}}{0.35}
    \figuresub[半空间]{\includegraphics{Halfspace.fig.pdf}}{0.35}
\end{Figure}

\subsection{多面体}

\begin{BoxDefinition}[多面体]
    多面体（Polyhedra）定义为
    \begin{Equation}[多面体]
        \polyhedra=\qty{\vb{x}\mid\vb{A}\vb{x}\preceq\vb{b},\ \vb{C}\vb{x}=\vb{d}}
    \end{Equation}
\end{BoxDefinition}

其中符号$\preceq$代表对向量逐元素的$\leq$。严格来说，这里的多面体其实是超多面体，若以$\R^3$为例，$\vb{A}\vb{x}\preceq\vb{b}$的意义是一系列包围出多面体的平面，$\vb{C}\vb{x}=\vb{d}$是一个额外的平面从多面体中切割出一个多边形。这里$\vb{C}$和$\vb{d}$也可以是空的，可以不对$\R^n$下得到的$n$维超多面体降维。

\subsection{球与椭球}

\begin{BoxDefinition}[球]
    球（Ball）定义为
    \begin{Equation}[球1]
        B(\vb{x}_c,r)=\qty{\vb{x}\mid(\vb{x}-\vb{x}_c)^2r^{-2}\leq 1}
    \end{Equation}
    也可以用参数方程的方式定义为
    \begin{Equation}[球2]
        B(\vb{x}_c,r)=\qty{\vb{x}+r\vb{u}\mid\norm{u}_2\leq 1}
    \end{Equation}
\end{BoxDefinition}

\begin{BoxDefinition}[椭球]
    椭球（Ellipsoid）定义为
    \begin{Equation}[椭球1]
        E(\vb{x}_c,\vb{P})=\qty{\vb{x}\mid(\vb{x}-\vb{x}_c)^TP^{-1}(\vb{x}-\vb{x}_c)\leq 1}
    \end{Equation}
    也可以用参数方程的方式定义为
    \begin{Equation}[椭球2]
        E(\vb{x}_c,\vb{P})=\qty{\vb{x}+\vb{A}\vb{u}\mid\norm{u}_2\leq 1}
    \end{Equation}
    其中$\vb{P}\in\Symm^n_{++}$是对称正定矩阵，而$\vb{A}=\vb{P}^{1/2}$。
\end{BoxDefinition}

$\vb{P}$的特征向量指向椭球轴的方向，$\vb{P}$的特征值的平方根$\sqrt{\lambda_i}$是椭球的半轴长。规定$\vb{P}$是对称正定矩阵的原因是避免出现更一般的二次曲面。特别的，当$\vb{P}=r^2\vb{I}$时，椭球会退化为球。

\subsection{范数球与范数锥}

\begin{BoxDefinition}[范数]
    若一个运算$\norm{\cdot}$满足
    \begin{enumerate}
        \item 正定性：$\norm{\vb{x}}\geq 0$，$\norm{\vb{x}}=0$当且仅当$\vb{x}=0$。
        \item 齐次性：$\norm{t\vb{x}}=\abs{t}\norm{x}$，对于$t\in\R$。
        \item 三角不等式：$\norm{\vb{x}+\vb{y}}\leq\norm{x}+\norm{y}$
    \end{enumerate}
    那么就可以将这种运算称之为范数（Norm）。
\end{BoxDefinition}

范数概括了一类具有距离性质的运算，实践中比较常用的是$p$\,-范数
\begin{Equation}[p-范数]
    \norm{\vb{x}}_{p}=\qty(\abs{x_1}^p+\cdots\abs{x_n}^p)^{1/p}
\end{Equation}

特别的，$1$-范数是曼哈顿距离$\norm{x}_1=\abs{x}_1+\cdots\abs{x}_n$，$2$-范数是欧氏距离$\sqrt{x_1^2+\cdots+x_n^2}$。

\begin{BoxDefinition}[范数球]
    范数球（Norm Ball）定义为
    \begin{Equation}[范数球]
        \qty{\vb{x}\in\R^n\mid\norm{\vb{x}-\vb{x}_c}\leq r}
    \end{Equation}
\end{BoxDefinition}

\begin{BoxDefinition}[范数锥]
    范数锥（Norm Cone）定义为
    \begin{Equation}[范数锥]
        \qty{(\vb{x},t)\in\R^{n+1}\mid\norm{\vb{x}}\leq t}
    \end{Equation}
\end{BoxDefinition}

例如，$\R^2$下的$2$-范数球就是$\sqrt{x^2+y^2}\leq r$，$\R^2$下的$2$-范数锥就是$|x|\leq y$，请注意范数锥的定义中$t$也会占用一个空间维度：在$\R^2$的$(x,y)$下$t=y$，在$\R^3$的$(x,y,z)$下$t=z$。

\subsection{半正定锥}

\begin{BoxDefinition}[对称]
    $\Symm^n$表示由对称（Symmetric）的$n\times n$方阵构成的集合。
    \begin{Equation}[对称]
        \Symm^n=\qty{\vb{X}\in\Symm^n\mid\vb{X}=\vb{X}^T}
    \end{Equation}
\end{BoxDefinition}

当我们说一个矩阵是正定或半正定时，它至少是对称的。

\begin{BoxDefinition}[正定]
    $\Symm^n_{++}$表示由对称正定（Positive Definite）的$n\times n$方阵构成的集合
    \begin{Equation}[正定]
        \Symm^n_{++}=\qty{\vb{X}\in\Symm^n\mid\vb{X}\succ 0}
    \end{Equation}
    若$\vb{X}\in\Symm^n_{++}$，这等价于对于任意的$\vb{z}\in\R^n$
    \begin{Equation}[正定的性质]
        \vb{z}^T\vb{X}\vb{z}> 0
    \end{Equation}
\end{BoxDefinition}

\begin{BoxDefinition}[半正定]
    $\Symm^n_{+}$表示由对称半正定（Positive Semidefinite）的$n\times n$方阵构成的集合
    \begin{Equation}[半正定]
        \Symm^n_{+}=\qty{\vb{X}\in\Symm^n\mid\vb{X}\succeq 0}
    \end{Equation}
    若$\vb{X}\in\Symm^n_{+}$，这等价于对于任意的$\vb{z}\in\R^n$
    \begin{Equation}[半正定的性质]
        \vb{z}^T\vb{X}\vb{z}\geq 0
    \end{Equation}
\end{BoxDefinition}

有趣的是，半正定集$\vb{S}^n_{+}$是一个凸锥，也称为半正定锥（Positive Semidefinite Cone）。

请注意，这里的$\succ$和$\succeq$不能理解为逐元素比较！严格来说，这两个符号代表的是“向量不等性”或“矩阵不等性”，在向量下确实会退化到逐元素比较，但矩阵下必须用二次型来定义。

不妨考虑这样一个例子，在$n=2$下的半正定条件是
\begin{Equation}[半正定例子1]
    \vb{X}=\mqty[
        x & y \\
        y & z \\
    ]\in\Symm^2_{+}\quad\Longleftrightarrow\quad
    x\geq 0,\ z\geq 0,\ xz\geq y^2
\end{Equation}

那么，下面这个矩阵是满足半正定条件的，尽管存在负元素
\begin{Equation}[半正定例子2]
    \vb{X}=\mqty[
        2 & -1\\
        -1 & 3\\
    ] 
\end{Equation}